% GENERATED FILE. Edit canonical lessons, then run npm run generate:books.
% canonical-source-hash: fnv1a64:a44cf4254a0c8425
% canonical-lessons: UR-C03-aap-tum-tu, UR-C03-kya, UR-C03-aap-ka-naam-kya-hai, UR-C03-khushi-hui, UR-C03-practice

\chapter{Ask and Answer Names}
\label{ch:ur-ask-and-answer-names}
\hlchaptermodality{3}

\begin{chapteropening}
\textbf{By the end of this chapter:} \emph{I can ask someone's name with respectful āp, answer with my own, and close the introduction politely.}

A six-line first meeting: greet, ask āp kā nām kyā hai?, answer with your own name, and close with āp se mil kar khushī huī.
\end{chapteropening}

\section[āp / tum / tū]{\ur{آپ} / \ur{تم} / \ur{تو} — three relationships inside “you”}

\label{lesson:UR-C03-aap-tum-tu}



\begin{quote}
Say \textbf{nahī̃} once, then \textbf{merā nām ... hai} with your own name. To ask the other person's name, begin with the respectful form.
\end{quote}

\subsection*{You'll want to know first — three forms}
\noindent
\begin{tabularx}{\linewidth}{@{}>{\raggedright\arraybackslash}X>{\raggedright\arraybackslash}X>{\raggedright\arraybackslash}X@{}}
\toprule
\textbf{Urdu} & \textbf{read} & \textbf{use} \\
\midrule
\textbf{\ur{آپ}} & \emph{āp} & respectful: a new adult, elder, or anyone you wish to honor \\
\textbf{\ur{تم}} & \emph{tum} & familiar: friends and peers \\
\textbf{\ur{تو}} & \emph{tū} & intimate, or downward; unsafe as a beginner default \\
\bottomrule
\end{tabularx}

\textbf{\ur{آپ}} begins with \textbf{\ur{آ}}, an alif carrying the long \emph{ā} mark, then \textbf{\ur{پ}} \emph{p}. \textbf{\ur{تم}} and \textbf{\ur{تو}} share \textbf{\ur{ت}} \emph{t}. In \textbf{\ur{تو}}, \textbf{\ur{و}} carries long \emph{ū}.

\begin{grammarlens}[title={Grammar Lens: choose the relationship first}]
English flattened these choices into one \emph{you}. Urdu keeps three. Use \textbf{āp} in a first meeting. Move to \textbf{tum} when the relationship supports familiarity. Save \textbf{tū} for relationships where you already know it is welcome; used badly, it can sound belittling.
\end{grammarlens}

\begin{cousinweb}
\textbf{Tum/tū} belong to the old Indo-European familiar-“you” family: Sanskrit \emph{tvam}, English \textbf{thou}, and Latin \emph{tū} are relatives. \textbf{Āp} followed another route: a word connected with “self/oneself” became honorific address.
\end{cousinweb}

\subsection*{Your turn}
Say these aloud:

\begin{itemize}
\raggedright
  \item a new adult — \textbf{āp}
  \item a friend — \textbf{tum}
  \item the old family — \textbf{tum, tvam, thou}
\end{itemize}

\begin{tcolorbox}[breakable,colback=teal!4,colframe=teal!35!black,title={Before you move on}]
Which form is safest when unsure? (\textbf{Āp}.) Which familiar forms share history with English \emph{thou}? (\textbf{Tum/tū}.)

Source: \href{https://openbooks.lib.msu.edu/urdu/chapter/2-8/}{\emph{Basic Urdu}: Cultural Notes}.
\end{tcolorbox}
\section[kyā]{\ur{کیا} — what}

\label{lesson:UR-C03-kya}



\begin{quote}
Which “you” opens a respectful first meeting? (\textbf{\ur{آپ}}, \emph{āp}.) Now add the question word.
\end{quote}

\subsection*{You'll want to know first — \ur{کیا}}
\begin{quote}
\textbf{\ur{کیا}} — \emph{kyā} — \textbf{what}
\end{quote}

Read right to left: \textbf{\ur{ک}} \emph{k}, \textbf{\ur{ی}} supplying the \emph{y} glide here, then \textbf{\ur{ا}} long \emph{ā}. The same \textbf{\ur{ی}} carried long \emph{ī} in \textbf{\ur{جی}} and \textbf{\ur{نہیں}}. Its job depends on the word; here the learned shape is \emph{kyā}.

\begin{cousinweb}
Urdu \textbf{kyā} and Hindi \textbf{kyā} are the same inherited Indo-Aryan word in different scripts. Their history reaches Sanskrit \emph{kim} and the old Indo-European \emph{k/kw-} question family behind English \textbf{wh-} words such as \emph{what, who,} and \emph{which}. The bridge helps memory; Urdu \textbf{\ur{کیا}} remains the form you must read here.
\end{cousinweb}

\subsection*{Your turn}
\begin{itemize}
\raggedright
  \item \emph{Say it:} \textbf{kyā} — what
  \item \emph{Read it:} \textbf{\ur{کیا}} as \textbf{k + yā}
  \item \emph{Say it:} the family — \textbf{kyā, what, who, which}
\end{itemize}

\begin{tcolorbox}[breakable,colback=teal!4,colframe=teal!35!black,title={Before you move on}]
Which letter acts as \emph{y} here but as long \emph{ī} in \textbf{\ur{جی}}? (\textbf{\ur{ی}}.) What does \textbf{\ur{کیا}} mean? (\textbf{What}.)
\end{tcolorbox}
\section[āp kā nām kyā hai?]{\ur{آپ} \ur{کا} \ur{نام} \ur{کیا} \ur{ہے؟} — ask a name respectfully}

\label{lesson:UR-C03-aap-ka-naam-kya-hai}



\begin{quote}
Recall \textbf{āp} “you,” \textbf{kyā} “what,” and the answer \textbf{merā nām ... hai}. The question reuses all three.
\end{quote}

\subsection*{The exchange}
\begin{quote}
\textbf{\ur{آپ} \ur{کا} \ur{نام} \ur{کیا} \ur{ہے؟}} — \emph{āp kā nām kyā hai?} — \textbf{What is your name?}
\end{quote}

Build it in spoken order:

\begin{quote}
you + \textbf{kā} + name + what + is?
\end{quote}

\textbf{Āp kā} means “your” in this frame. The possessive form agrees with the thing owned: \textbf{nām} is masculine singular, so use \textbf{kā}. Keep that one pairing; the full \emph{kā/kī/ke} system can wait. As in your answer, \textbf{hai} stays at the end.

\subsection*{Script: the question mark}
Urdu ends the question with \textbf{\ur{؟}}, curved for a right-to-left line. Begin at \textbf{\ur{آپ}} on the right and reach the punctuation at the far left.

\subsection*{Your turn}
\begin{itemize}
\raggedright
  \item \emph{Say it:} \textbf{āp kā nām kyā hai?}
  \item \emph{Say it:} ask, then answer — \textbf{merā nām ... hai}
  \item \emph{Point to:} to \textbf{\ur{؟}} at the end
\end{itemize}

\begin{tcolorbox}[breakable,colback=teal!4,colframe=teal!35!black,title={Before you move on}]
Ask once without romanization. Where does \textbf{hai} remain? (At the end.)

Source: \href{https://openbooks.lib.msu.edu/urdu/chapter/2-2/}{\emph{Basic Urdu}: Greetings and Introductions}.
\end{tcolorbox}
\section[āp se mil kar khushī huī]{\ur{آپ} \ur{سے} \ur{مل} \ur{کر} \ur{خوشی} \ur{ہوئی} — pleased to meet you}

\label{lesson:UR-C03-khushi-hui}



\begin{quote}
Say the opening ladder once: \textbf{shukriyā — jī hā̃ — nahī̃}. Point to the three shapes you traced before the joined thank-you word. Then ask \textbf{āp kā nām kyā hai?} and answer \textbf{merā nām ... hai}.
\end{quote}

\subsection*{The exchange}
\begin{quote}
\textbf{\ur{آپ} \ur{سے} \ur{مل} \ur{کر} \ur{خوشی} \ur{ہوئی۔}} \emph{āp se mil kar khushī huī} \textbf{Pleased to meet you.}
\end{quote}

For now, hold three meaning-sized chunks:

\begin{itemize}
\raggedright
  \item \textbf{āp se} — with you
  \item \textbf{mil kar} — having met
  \item \textbf{khushī huī} — happiness happened
\end{itemize}

That yields the memorable literal picture: \textbf{“Having met you, happiness happened.”} The phrase is one social move; later lessons can generalize \emph{se}, \emph{-kar}, and past agreement one at a time.

\begin{cousinweb}
\textbf{Khushī} “happiness” comes from Persian \textbf{khosh/khush} “pleasant, happy.” \textbf{Mil} “meet” and \textbf{huī} “happened” belong to Urdu's inherited Indo-Aryan grammar. The sentence therefore joins Urdu's two major historical layers without becoming any less ordinary Urdu.
\end{cousinweb}

\subsection*{Your turn}
Say these aloud:

\begin{itemize}
\raggedright
  \item \textbf{āp se | mil kar | khushī huī} in three beats
  \item the literal picture — “having met you, happiness happened”
  \item which word came through Persian — \textbf{khushī}
\end{itemize}

\begin{tcolorbox}[breakable,colback=teal!4,colframe=teal!35!black,title={Before you move on}]
Say the response once as a whole. Which chunk means “happiness happened”? (\textbf{Khushī huī}.)

Source: \href{https://openbooks.lib.msu.edu/urdu/chapter/2-2/}{\emph{Basic Urdu}: Greetings and Introductions}.
\end{tcolorbox}
\section[Practice]{Practice — a complete Urdu name exchange}

\label{lesson:UR-C03-practice}



\begin{quote}
Recall the four moves without adding anything: greet, ask, answer, acknowledge.
\end{quote}

\subsection*{The exchange}
\begin{quote}
A: \textbf{\ur{سلام}!} — \emph{Salām!} B: \textbf{\ur{سلام}!} — \emph{Salām!} A: \textbf{\ur{آپ} \ur{کا} \ur{نام} \ur{کیا} \ur{ہے؟}} — \emph{Āp kā nām kyā hai?} B: \textbf{\ur{میرا} \ur{نام} \ur{سارا} \ur{ہے۔}} — \emph{Merā nām Sārā hai.} A: \textbf{\ur{آپ} \ur{سے} \ur{مل} \ur{کر} \ur{خوشی} \ur{ہوئی۔}} — \emph{Āp se mil kar khushī huī.} B: \textbf{\ur{شکریہ۔}} — \emph{Shukriyā.}
\end{quote}

Every line is already known. \textbf{Āp} keeps the meeting respectful; \textbf{hai} stays at the end of both name sentences; \textbf{\ur{؟}} ends the question at the line's left.

\subsection*{Your turn}
Say these aloud:

\begin{itemize}
\raggedright
  \item both voices, using your own name
  \item the exchange again without romanization
  \item only the three name moves — question, answer, response
\end{itemize}

\begin{tcolorbox}[breakable,colback=teal!4,colframe=teal!35!black,title={Before you move on}]
Run the meeting once from \textbf{salām} to \textbf{shukriyā}. If a line stalls, review only that micro-lesson.
\end{tcolorbox}
